% !TEX program = lualatex
\documentclass[11pt,a4paper]{article}
\usepackage{icat-common}

%-------------------------------------------------
%  独自コマンド・設定
%-------------------------------------------------

% セクション名をラテン語形式に変更

%-------------------------------------------------
\begin{document}

\title{直観主義四句分別における真理保存性}
\author{千葉将臣}
\date{2026-04-19}
\maketitle

\begin{abstract}
本稿は、真理保存性（Truth Preservation）の基盤を静的な自己同一性から動的なフラクタル性へと転回させ、それが体系内部でいかに記述可能であるかを論証するものである。二値論理が抱えるメタ理論的限界を、双方向演繹定理と構成プロセスによって克服する機序を明らかにする。
\end{abstract}

\section*{序言}
二値論理における真理保存性は、真理値の静的な自己同一性に依拠し、体系外のメタ理論（健全性定理等）としてのみ扱われてきた。これに対し、直観主義四句分別においては、真理保存性そのものが「フラクタル構造の維持」として体系内部のプロセスへと取り込まれる。これにより、勝義諦（極限）と世俗諦（近似）の連続性が論理的に担保される。

\section{定義}

\begin{description}[style=multiline, leftmargin=4.5cm]
  \item[定義 4.1] \textbf{自己同一性的保存}：真理値を静的な「状態（0または1）」として固定し、推論をその状態の不変的な変換として扱う機序。
  \item[定義 4.2] \textbf{フラクタル的保存}：真理を「構成プロセス」として定義し、そのプロセスがスケールを越えてフラクタル則（自己相似性等）を維持することを指す。
  \item[定義 4.3] \textbf{体系内記述可能性}：論理体系の妥当性や真理保存性が、メタ言語を介さず、体系内部の対象（証明、近似列、極限状態）として表現・演算できる性質。
\end{description}

\section{公理}

\begin{description}[style=multiline, leftmargin=4.5cm]
  \item[公理 4.1] 二値論理における真理保存性は、体系外の意味論（モデル）に依存する。
  \item[公理 4.2] 直観主義四句分別における真理保存性は、構成プロセスの幾何学的形状（フラクタル）に依存する。
\end{description}

\section{定理}

\begin{description}[style=multiline, leftmargin=4.5cm]
  \item[定理 4.1] \textbf{真理保存性のフラクタル的基盤}\\
  直観主義四句分別における推論の妥当性は、証明の構成プロセスがフラクタル則（定理 2.1--2.4）を満たすことによって担保される。真理保存性は「自己同一性」ではなく「構造の自己相似的な反復」に帰着する。
  
  \item[定理 4.2] \textbf{真理保存性の体系内統合}\\
  真理が構成可能性として定義される（定義 3.1）ことにより、ある証明から別の証明への遷移、および近似列の極限状態（定理 3.2）への収束は、体系内部の対象として記述可能となる。ゆえに真理保存性は体系内の演算として扱われる。
\end{description}

\section{補足}
二値論理が真理をメタ理論へと追いやったのは、高次元の「ない（非）」を低次元の「ある（是）」の反転として同次元に閉じ込めた結果である。これに対し、直観主義四句分別の体系では、近似（世俗諦）が極限（勝義諦）へと収束するプロセスそのものを真理保存と見做す。これは、真理を「点（状態）」ではなく「線（プロセス・螺旋）」として捉える転回である。

\section{系}
したがって、体系内で非収束の余剰を検知・抑止する機能は、このフラクタル的真理保存性の帰結として自然に要請される。これは正当性（Justification）とは、プロセスがフラクタルな縮小写像として安定していることの別の表現に他ならない。

\end{document}
